\documentclass{article}
\usepackage[a3paper,landscape,margin=15mm]{geometry}
\usepackage[UTF8,fontset=fandol]{ctex}
\usepackage[hidelinks]{hyperref}
\usepackage{./tex/mindchart}
\pagestyle{empty}
\begin{document}
\section*{Publication layout / 技术方案}
\noindent
\begin{mindchart}[layout=balanced,fit_width=\linewidth,fit_height=19cm,avoid_overlaps=true,max_width=3cm,root_style={font=\Large,fill=blue,text_color=white},branch_style={font=\normalsize,shape=card,pad_x=10},leaf_style={font=\small,shape=rect,fill=white}]
# 系统设计, name=root
  - 数据处理与质量验证, name=data, side=left, tag=v2
    - 长标题：输入数据的清洗、归一化与完整性检查, name=clean
    - 特征变换 $z=\frac{x-\mu}{\sigma}$, name=feature
  - 模型训练, name=model, side=right
    - 编码器, name=enc
      - 多层特征融合, name=fusion
        - 损失函数 $\mathcal{L}(y,\hat{y})$
    - 分类器, name=classifier
  - 评估与发布, name=release, side=right, href=https://example.com
    - 验证集评估, status=Done
\summary{enc}{classifier}{模型组件}
\summary{enc}{classifier}{第二组概要}
\boundary{data}{数据质量要求}
\boundary{model}{训练模块}
\note[name=assumption,side=bottom,max_width=4cm,node_ref=clean]{data}{假设条件：输入已完成脱敏。点击可跳转至清洗节点。}
\rel{feature}{enc}{特征输入}
\rel{fusion}{release}{模型交付}
\end{mindchart}
\end{document}
